In this appendix, we provide the technical details of the proof of \Cref{thm:synthetic-Adams}, as well as a computation of the $\HFt$-synthetic bigraded homotopy groups in the Toda range. The computation of synthetic homotopy groups highlights many of the subtleties within the statement of \Cref{thm:synthetic-Adams}. We have tried to make this appendix as self-contained as possible. % and except for a couple uses of \Cref{lemm:strong-etau-complete} in later sections we have succeeded.
Understanding the techniques introduced in this appendix is not necessary in order to read the remainder of the paper. For convenience, we recall the statement of \Cref{thm:synthetic-Adams}.

\begin{thm}[\Cref{thm:synthetic-Adams}]
  Let $X$ denote an $E$-nilpotent complete spectrum with strongly convergent $E$-based Adams spectral sequence.
  Then we have the following description of the bigraded homotopy groups of $\nu X$.

  Let $x$ denote a class in topological degree $k$ and filtration $s$ of the $\mathrm{E}_2$-page of the $E$-based Adams spectral sequence for $X$. The following are equivalent:
  \begin{itemize}
  \item[(1a)] Each of the differentials $d_2$,...,$d_r$ vanish on $x$.
  \item[(1b)] $x$, viewed as an element of $\pi_{k,k+s}(C\tau \otimes \nu X)$, lifts to $\pi_{k,k+s}(C\tau^r \otimes \nu X)$.
  \item[(1c)] $x$ admits a lift to $\pi_{k,k+s}(C\tau^r \otimes \nu X)$ whose image under the $\tau$-Bockstein
  $$ C\tau^r \otimes \nu X \to \Sigma^{1,-r} C\tau \otimes \nu X $$
  is equal to $-d_{r+1}(x)$.
  \end{itemize}

  If we moreover assume that $x$ is a permanent cycle, then there exists a (not necessarily unique) lift of $x$ along the map
  $\pi_{k,k+s}(\nu X) \to \pi_{k,k+s}(C\tau \otimes \nu X)$.  For any such lift, $\wt{x}$, the following statements are true:
  \begin{itemize}
  \item[(2a)] If $x$ survives to the $\mathrm{E}_{r+1}$-page, then $\tau^{r-1} \wt{x} \neq 0$.
  \item[(2b)] If $x$ survives to the $\mathrm{E}_\infty$-page, then the image of $\wt{x}$ in $\pi_{k} (X)$ is of $E$-Adams filtration $s$ and detected by $x$ in the $E$-based Adams spectral sequence.
  \end{itemize}
  Furthermore, there always exists a choice of lift $\wt{x}$ satisfying additional properties:
  \begin{itemize}
  \item[(3a)] If $x$ is the target of a $d_{r+1}$-differential, then we may choose $\wt{x}$ so that $\tau^r \wt{x} = 0$.
  \item[(3b)] If $x$ survives to the $\mathrm{E}_\infty$-page, and $\alpha \in \pi_k X$ is detected by $x$, then we may choose $\wt{x}$ so that
    $\tau^{-1} \wt{x} = \alpha$. In this case we will often write $\wt{\alpha}$ for $\wt{x}$.
  \end{itemize}
  Finally, the following generation statement holds:
  \begin{itemize}
  \item[(4)] Fix any collection of $\wt{x}$ (not necessarily chosen according to (3)) such that the $x$ span the permanent cycles in topological degree $k$.  Then the $\tau$-adic completion of the $\Z[\tau]$-submodule of $\pi_{k,*}(\nu X)$ generated by those $\wt{x}$ is equal to $\pi_{k,*} (\nu X)$.
  \end{itemize}
\end{thm}

\begin{rmk}
  As a foreward to the proof, we provide some commentary on the provenance of \Cref{thm:synthetic-Adams}.
  Through the equivalence between $(\Syn_{\BP}^{\mathrm{ev}})_p$ and $\mathcal{SH}(\mathbb{C})_p^{\mathrm{cell}}$ (see \cite[Theorem 1.4]{Pstragowski}) this theorem specializes to provide a translation between the $p$-complete, bigraded motivic stable stems over $\mathbb{C}$ and the Adams--Novikov spectral sequence.
  In fact, the proof we give was directly inspired by the literature on motivic stable stems and their connection to the Adams--Novikov spectral sequence.

  The core argument originates in \cite[Lemma 15]{HKO} where is it observed that at the prime $2$ the differentials in the motivic Adams--Novikov spectral sequence can be formally deduced from those in the classical Adams--Novikov spectral sequence.\footnote{See \cite[Proposition 5.6]{Stahn} for a version at odd primes.} The corresponding result in our setting is \Cref{thm:synth-ASS} where we identify the $\nu E$-based Adams spectral sequence for $\nu X$ in terms of the $E$-based Adams spectral sequence for $X$.

  Building on knowledge from extensive calculations of motivic stable stems, Isaksen then recognized that not only do these spectral sequences determine one another, but the motivic stable stems and Adams--Novikov spectral sequence (with its hidden extensions) contain essentially equivalent information (see the introduction to \cite[Chapter 6]{StableStems}). The remainder of the proof of \Cref{thm:synthetic-Adams} involves formalizing these ideas and dealing with questions of convergence.
  \todo{I'm not too happy with this remark, examine again with fresh eyes.}
\end{rmk}

%%% Local Variables:
%%% mode: latex
%%% eval: (visual-line-mode t)
%%% eval: (auto-fill-mode 0)
%%% End:
